% ============================================================
\section{Implementation Gap and Stub Audit}
\label{sec:implementation_gap_audit}
% ============================================================

This chapter records the current code-backed boundary between implemented
runtime behavior, transparent approximations, resolved gaps, and remaining
bounded model scope.  It is intentionally separate from the mathematical model
chapters: a formula in this notebook is not a claim that every richer physical
variant is solved end-to-end in production code.  Each status item below is tied
to a concrete API, source path, and regression test rather than to narrative
intent alone.

\subsection{Audit Method}

The audit covered \verb|src/|, \verb|include/|, \verb|tests/|, the
technical-notebook \verb|sections/| directory, and the sibling external
\verb|MIPSolvers| checkout used by CMake.  The review combined source inspection
with searches for high-signal markers such as \verb|stub|, \verb|unimplemented|,
\verb|placeholder|, \verb|fallback|, \verb|simplified|, \verb|approximate|,
\verb|dummy|, and \verb|unsupported|.  Marker counts are not used as defect
counts: many occurrences are legitimate solver fallback paths, explicit
input-domain checks, or tests that verify a declared boundary.

The latest pass specifically audited external solver reproducibility, native
B\&C quiet-mode behavior, hybrid topology fault addressing, topology load
accounting, legacy ONR branch-ID semantics, three-stage reliability exactness,
graph contraction domain maps, and DCOPF branch-dual validity.  A second claim
audit then re-read this chapter against the code and tests.  It found no stale
closed-gap statement, but it did expose one additional quiet-mode leak in the
strict resilience MIP path; that leak is now gated in the external solver and
covered by file-descriptor-level output capture.

\subsection{High-Risk Items Closed}

\begin{enumerate}
  \item \textbf{External MIPSolvers reproducibility.}  The external dependency
        is expected at pinned commit
        \verb|f0916c99f7f1364b825cabb58b50c3cc3d20eddd|.  CMake treats any
        external \verb|MIPSolvers| HEAD mismatch as fatal, and now also treats a
        dirty external checkout as fatal for Release, RelWithDebInfo,
        MinSizeRel, and CI configurations.  The local override
        \verb|HACDCDSS_SKIP_MIPSOLVERS_DIRTY_CHECK| is rejected in those strict
        modes, so a reproducible Release artifact cannot silently include
        unpinned or uncommitted external solver edits.

  \item \textbf{Distribution power-flow public API.}  The legacy distribution
        power-flow compatibility header and public single-phase stub/BFS-style
        API are removed.  The public distribution feeder surface is the
        three-phase NR contract in
        \verb|include/hacdcpf/power_flow/distribution_power_flow.hpp|:
        \verb|solve_three_phase_nr| and phase-domain Y-bus helpers.  Searches
        find no live \verb|solve_distribution_pf|, \verb|DistributionPFOptions|,
        or \verb|DistributionPFResult| code symbols.

  \item \textbf{Hybrid FMEA physical evaluation.}  Hybrid FMEA no longer scores
        DC/VSC systems through an AC-only DCOPF boundary.  When DC or converter
        components exist, \verb|run_distribution_fmea| uses a linear hybrid
        AC/DC network LP that co-optimizes AC load shedding, DC load shedding,
        AC/DC sources, DC branches, VSC active-power transfers, and DC/DC
        converter transfers.  Reports set \verb|model_scope| to
        \verb|hybrid-acdc-network-lp| and mark DC-load curtailment plus VSC/DC
                        power flow as included.  The same validity contract now explicitly
                        leaves \verb|ac_voltage_reactive_feasibility_certified=false| because
                        nonlinear AC voltage/reactive checks are outside this LP evaluator.

  \item \textbf{AC OPF fallback for hybrid cases.}  Hybrid AC/DC cases are
        detected before fallback selection.  If the nonlinear hybrid path fails,
        \verb|solve_ac_opf| suppresses the AC-only economic-dispatch plus AC-PF
        fallback and attaches an explicit diagnostic that the case contains
        DC/VSC components.  The fast AC-only route remains available for AC-only
        cases.
\end{enumerate}

\subsection{Other Resolved or Explicitly Bounded Findings}

\begin{enumerate}
  \item \textbf{Native B\&C quiet mode.}  The external native branch-and-cut
        implementation now gates root-phase timing, conformance, clique,
        propagation, reduced-cost lurking, implied-bound, proof, objective, and
        queue diagnostics behind
        \verb|BCOptions::verbose| or explicit diagnostic environment variables
        such as \verb|MIPSOLVERS_BC_CONF| and \verb|MIPSOLVERS_BC_TIMELINE|.
        File-descriptor-level stdout/stderr capture tests in topology
        reconfiguration, legacy ONR, and the strict resilience MIP entry point
        verify that \verb|verbose=false| is actually quiet.

  \item \textbf{Three-stage reliability exactness.}
        \verb|src/reliability/three_stage_reliability.cpp| is a native C++
        evaluator.  It solves finite-source AC LinDistFlow restoration with
        explicit source bounds, load-shed variables, branch limits, voltage
        bounds, strict radial forest constraints, healthy normally-closed edge
        coupling, and Stage 2 switch-operation limits.  True MILPs use HiGHS
        first and native B\&C as fallback; per-stage status and MIP gaps are
        surfaced.  No Julia or Gurobi runtime is required.

        Result field \verb|ok| now means verified exactness inside the
        implemented physical scope, not merely successful metric production.
        Pure AC cases can return \verb|ok=true|.  Hybrid DC/VSC/DCDC/SOP cases
        return metrics with \verb|ok=false|, nonempty \verb|error|, and model
        scope \verb|ac-lindistflow-milp+dc-connectivity-fallback|.  Regression
        tests protect pure-AC exactness, hybrid non-exactness, finite source
        capacity, AC bus load de-duplication, the no-free-open rule for healthy
        normally-closed lines, and VSC-capacity-limited AC support to DC load.

  \item \textbf{Topology reconfiguration hybrid addressing.}
        \verb|src/network_reconfiguration/topology_reconfiguration.cpp| tries
        HiGHS when available and otherwise calls the native B\&C fallback.
        The result exposes \verb|solver_backend|, \verb|solver_status|,
        \verb|solver_mip_gap|, and \verb|proven_optimal|; \verb|optimal| mirrors
        that certificate.  Switchable-edge input accepts structured
        \verb|BranchRef| values via \verb|switchable_branches|; fault input now
        also accepts structured \verb|faulted_branches|.  The legacy
        \verb|switchable_branch_ids| and \verb|line_failures| fields are AC-only
        compatibility fields.

  \item \textbf{Topology load accounting.}  Fault reconfiguration adds explicit
        \verb|ACSystem::loads| demand, direct \verb|ACBus::pd_mw| and
        \verb|ACBus::qd_mvar| demand, and direct \verb|DCBus::pd_mw| demand
        instead of treating any one record type as a replacement for the others.

  \item \textbf{Legacy ONR branch identity.}  The legacy AC ONR API documents
        \verb|ONROptions::switchable_branch_ids| as stable \verb|ACBranch::index|
        values, not vector positions.  Non-candidate branches are fixed to
        their current \verb|in_service| state, and fixed \verb|alpha| variables
        are continuous constants omitted from the native B\&C binary set.

  \item \textbf{Resilience MIP entry point.}
        \verb|run_distribution_resilience_mip_assessment| builds an independent
        multi-period LinDistFlow-style restoration MILP.  It does not delegate
        to the heuristic assessment.  Result statistics expose
        \verb|model_scope|, validity flags, solver name/status, objective, gap,
        and constraint counts.  The MIP scope remains AC-only LinDistFlow and
        reports that scope explicitly through \verb|model_stats.model_scope|,
        \verb|validity.dc_network_modelled=false|, and
        \verb|validity.vsc_dispatch_modelled=false|.

  \item \textbf{DC OPF load-shedding audit.}  Direct isolated-island shedding is
        written to both \verb|total_load_shedding_mw| and the per-bus
        \verb|load_shedding_mw| vector.  The topology pre-prune removes loads,
        generators, fixed injections, storage, and AC branches attached to dead
        islands before LP construction.  The independent feasibility check
        includes slack-bus balance.

  \item \textbf{DC OPF JSON contract.}  DCOPF result JSON helpers
        serialize/restore solver audit fields, including \verb|solver_chain|
        and \verb|objective_model|.

  \item \textbf{DC OPF LMPs and branch duals.}  Nodal LMP extraction is
        implemented through a supporting simplex LP when \verb|compute_lmp=true|.
        Branch shadow-price arrays are intentionally flagged with
        \verb|branch_mu_valid=false| until bounded-variable congestion duals are
        recovered robustly.  The LMP regression asserts this uncertified
        branch-dual boundary explicitly.

  \item \textbf{Graph contraction and recovery maps.}  Hybrid graph contraction
        results expose domain-qualified maps such as \verb|ac_bus_to_super| and
        \verb|dc_bus_to_super|.  Recovery helpers can use the structured result
        directly so same numeric AC/DC bus IDs are not collapsed through a
        legacy combined map.

  \item \textbf{JPC rich-component round-trip.}  \verb|test_io_json| includes
        rich-component JPC round-trip coverage and asserts the populated
        component collections are preserved.
\end{enumerate}

\subsection{Closed Model-Scope Contracts}

The second pass treats the following items as closed contracts rather than open
implementation gaps.  They are not silently solved by approximation and they are
not advertised as richer physics than the code enforces; each boundary is
visible through a result field and protected by regression tests.

\begin{enumerate}
  \item Hybrid FMEA is a linear restoration/curtailment LP.  It includes DC
        loads, DC network edges, VSC active transfers, DC/DC transfers, and
        source limits, but it does not certify nonlinear AC voltage/reactive
        feasibility for every restoration state.  Callers can check
        \verb|FMEAResult::validity.ac_voltage_reactive_feasibility_certified|,
        which remains false for this evaluator.

  \item The resilience MIP entry point remains an AC-only LinDistFlow
        restoration MILP and reports that scope explicitly through
        \verb|model_stats.model_scope=ac-only-lindistflow| plus false DC/VSC
        validity flags.

  \item Three-stage reliability is exact only for the pure-AC implemented
        restoration MILP scope.  Hybrid DC/VSC/DCDC/SOP cases report useful
        connectivity/capacity metrics but clear \verb|ok| and set validity
        flags false for DC power flow and whole-system restoration MILP
        coverage.

  \item DC OPF branch shadow-price arrays remain invalid unless
        \verb|branch_mu_valid=true|; bounded-variable branch-dual extraction is
        still a deliberately exposed boundary.  Nodal LMPs remain available via
        the supporting simplex LP when \verb|compute_lmp=true|.
\end{enumerate}

\subsection{Verification Snapshot}

After the implementation pass, the following commands were run from the
repository root using the \verb|build| Debug tree.  A Release configure was also
attempted specifically to prove that external \verb|MIPSolvers| reproducibility
is enforced.  The latest rerun fails before dirty-tree checking because the
local external checkout HEAD is not the pinned commit; therefore the full
Release CTest suite is intentionally blocked until the external solver tree is
checked out at the recorded pin and cleaned or the pin is intentionally updated.

\begin{verbatim}
cmake --build build -j4
./build/tests/test_three_stage_reliability "[reliability][three_stage][integration][milp]"
./build/tests/test_three_stage_reliability "[reliability][three_stage][regression]"
./build/tests/test_topology_crossval "[topology][reconfiguration][regression]"
./build/tests/test_topology_crossval "[topology][onr][regression]"
./build/tests/test_resilience_assessment "Resilience: MIP solver returns feasible on zero-fault system"
./build/tests/test_resilience_assessment "[fmea][reliability][regression]"
./build/tests/test_acopf_dcopf_crossval "[integration][opf][dcopf][lmp]"
./build/tests/test_graph "[graph][contraction]"
./build/tests/test_graph "[graph][recovery]"
./build/tests/test_graph "[graph][collision]"
cmake -S . -B /tmp/hacdcpf_release_dirty_check -DCMAKE_BUILD_TYPE=Release
ctest --test-dir build --output-on-failure
\end{verbatim}

Focused regression outcomes include: topology reconfiguration \textbf{28
assertions in 7 cases}, legacy ONR \textbf{6 assertions in 1 case}, strict
resilience MIP scope/quiet-mode \textbf{10 assertions in 1 case}, FMEA hybrid
scope regressions \textbf{16 assertions in 2 cases}, three-stage regressions
\textbf{24 assertions in 5 cases}, three-stage integration \textbf{1134
assertions in 3 cases}, graph contraction / recovery / collision \textbf{23 /
27 / 28 assertions}, and DCOPF LMP \textbf{41 assertions in 3 cases}.  The
Release configure fails as expected with a fatal reproducibility message; in the
latest rerun, CMake stopped at the external \verb|MIPSolvers| HEAD pin mismatch
before reaching the dirty-tree check.  The full Debug CTest result is
\textbf{432 registered tests}: \textbf{430 active tests passed}, \textbf{0
failed}, and \textbf{2 expected MATPOWER parser skips} (\verb|case533mt_hi|,
\verb|case533mt_lo|).

\subsection{Audit Conclusion}

The repository is not an empty shell: PF, OPF, graph, I/O, carbon, time-series,
resilience, FMEA, topology reconfiguration, and three-stage reliability paths
are implemented and covered by the current 432-entry CTest suite.  The three
previous high-risk audit items in this chapter are now closed in code and in
tests: the legacy distribution power-flow public API is removed, hybrid FMEA
optimizes DC networks and VSC transfers, and hybrid AC OPF can no longer accept
an AC-only fallback as a converged hybrid result.  This pass also closed the
newly audited logic holes in external MIPSolvers reproducibility, native B\&C
quiet mode, three-stage reliability exactness, topology reconfiguration, legacy
ONR branch identity, graph domain-map handling, DCOPF branch-dual validity, and
FMEA DC load scaling.  The second-pass review found no remaining unsupported
claim in this chapter; the items that would otherwise read as gaps are now
explicit model-scope contracts surfaced through result fields, validity flags,
diagnostics, and regression tests.